                              IMC 2012, Blagoevgrad, Bulgaria
                                           Day 2, July 29, 2012

Problem 1. Consider a polynomial

                                   f (x) = x2012 + a2011 x2011 + . . . + a1 x + a0 .

Albert Einstein and Homer Simpson are playing the following game. In turn, they choose one of the
coefficients a0 , . . . , a2011 and assign a real value to it. Albert has the first move. Once a value is assigned
to a coefficient, it cannot be changed any more. The game ends after all the coefficients have been
assigned values.
    Homer’s goal is to make f (x) divisible by a fixed polynomial m(x) and Albert’s goal is to prevent
this.

 (a) Which of the players has a winning strategy if m(x) = x − 2012?

 (b) Which of the players has a winning strategy if m(x) = x2 + 1?

                                            (Proposed by Fedor Duzhin, Nanyang Technological University)
Solution. We show that Homer has a winning strategy in both part (a) and part (b).
(a) Notice that the last move is Homer’s, and only the last move matters. Homer wins if and only if
f (2012) = 0, i.e.

                    20122012 + a2011 20122011 + . . . + ak 2012k + . . . + a1 2012 + a0 = 0.                           (1)

Suppose that all of the coefficients except for ak have been assigned values. Then Homer’s goal is to
establish (1) which is a linear equation on ak . Clearly, it has a solution and hence Homer can win.
(b) Define the polynomials

  g(y) = a0 + a2 y + a4 y 2 + . . . + a2010 y 1005 + y 1006       and h(y) = a1 + a3 y + a5 y 2 + . . . + a2011 y 1005 ,

so f (x) = g(x2 ) + h(x2 ) · x. Homer wins if he can achieve that g(y) and h(y) are divisible by y + 1, i.e.
g(−1) = h(−1) = 0.
    Notice that both g(y) and h(y) have an even number of undetermined coefficients in the beginning
of the game. A possible strategy for Homer is to follow Albert: whenever Albert assigns a value to a
coefficient in g or h, in the next move Homer chooses the value for a coefficient in the same polynomial.
This way Homer defines the last coefficient in g and he also chooses the last coefficient in h. Similarly
to part (a), Homer can choose these two last coefficients in such a way that both g(−1) = 0 and
h(−1) = 0 hold.

Problem 2. Define the sequence a0 , a1 , . . . inductively by a0 = 1, a1 = 21 and

                                                      na2n
                                        an+1 =                          for n ≥ 1.
                                                 1 + (n + 1)an
                        ∞
                        X ak+1
Show that the series               converges and determine its value.
                        k=0
                              ak
                                                       (Proposed by Christophe Debry, KU Leuven, Belgium)




                                                              1
Solution. Observe that
                                  (1 + (k + 1)ak )ak+1   ak+1
                          kak =                        =      + (k + 1)ak+1                              for all k ≥ 1,
                                           ak             ak
and hence
         n                    n
         X ak+1           a1 X                         1
                      =     +    (kak − (k + 1)ak+1 ) = + 1 · a1 − (n + 1)an+1 = 1 − (n + 1)an+1                                     (1)
         k=0
                 ak       a0 k=1                       2

for all n ≥ 0.
                          n
                          P                                                                                               ∞
                                                                                                                          P
                            ak+1                                                                                            ak+1
    By (1) we have               ak
                                        < 1. Since all terms are positive, this implies that the series                         ak
                                                                                                                                      is
                          k=0                                                                                             k=0
                                                 a
convergent. The sequence of terms, k+1
                                   ak
                                       must converge to zero. In particular, there is an index n0 such
     ak+1   1
that ak < 2 for n ≥ n0 . Then, by induction on n, we have an < 2Cn with some positive constant C.
From nan < Cn2n
                → 0 we get nan → 0, and therefore
                              ∞                       n                                  
                              X ak+1                  X ak+1
                                             = lim                   = lim 1 − (n + 1)an+1 = 1.
                                        ak      n→∞           ak          n→∞
                                k=0                   k=0


Remark. The inequality an ≤ 21n can be proved by a direct induction as well.

Problem 3. Is the set of positive integers n such that n! + 1 divides (2012n)! finite or infinite?
                                        (Proposed by Fedor Petrov, St. Petersburg State University)
Solution 1. Consider a positive integer n with n! + 1 (2012n)!. It is well-known that for arbitrary
nonnegative integers a1 , . . . , ak , the number (a1 + . . . + ak )! is divisible by a1 ! · . . . · ak !. (The number
of sequences consisting of a1 digits 1, . . . , ak digits k, is (aa11+...+a      k )!
                                                                         !·...·ak !
                                                                                      .) In particular, (n!)2012 divides
(2012n)!.
    Since n! + 1 is co-prime with (n!)2012 , their product (n! + 1)(n!)2012 also divides (2012n)!, and
therefore
                                           (n! + 1) · (n!)2012 ≤ (2012n)! .
                                          n
    By the known inequalities n+1       e
                                             < n! ≤ nn , we get
                       n 2013n
                                       < (n!)2013 < (n! + 1) · (n!)2012 ≤ (2012n)! < (2012n)2012n
                          e
                                                         n < 20122012 e2013 .

Therefore, there are only finitely many such integers n.
                                                 
                                              n+1 n
Remark. Instead of the estimate                e    < n!, we may apply the Multinomial theorem:
                                                                    X                N!
                                      (x1 + · · · + xℓ )N =                                       xk11 . . . xkℓ ℓ .
                                                                              k1 ! · . . . · kℓ !
                                                              k1 +...+kℓ =N

Applying this to N = 2012n, ℓ = 2012 and x1 = . . . = xℓ = 1,

                                          (2012n)!
                                                    < (1| + 1 +{z. . . + 1})2012n = 20122012n ,
                                           (n!)2012
                                                                   2012
                                                                    (2012n)!
                                               n! < n! + 1 ≤                  < 20122012n .
                                                                     (n!)2012



                                                                          2
    On the right-hand side we have a geometric progression which increases slower than the factorial
on the left-hand side, so this is true only for finitely many n.
Solution 2. Assume that n > 2012 is an integer with n! + 1 (2012n)!. Notice that all prime divisors
of n! + 1 are greater than n, and all prime divisors of (2012n)! are smaller than
                                                                                h 2012n.
                                                                                     i
    Consider a prime p with n < p < 2012n. Among 1, 2, . . . , 2012n there are 2012n
                                                                                   p
                                                                                       < 2012 numbers
divisible by p; by p2 > n2 > 2012n, none of them is divisible by p2 . Therefore, the exponent of p in the
prime factorization of (2012n)! is at most 2011. Hence,
                                                                Y
                            n! + 1 = gcd n! + 1, (2012n)! <            p2011 .
                                                                              n<p<2012p
                           Q
Applying the inequality         p < 4X ,
                          p≤X
                                                                  !2011
                           Y                         Y                               2011                  n
                                         2011
                  n! <               p          <             p           < 42012n           = 42012·2011        .   (2)
                         n<p<2012p                  p<2012n

Again, we have a factorial on the left-and side and a geometric progression on the right-hand side.

Problem 4. Let n ≥ 2 be an integer. Find all real numbers a such that there exist real numbers x1 ,
. . . , xn satisfying
                    x1 (1 − x2 ) = x2 (1 − x3 ) = . . . = xn−1 (1 − xn ) = xn (1 − x1 ) = a.                         (1)
                                                (Proposed by Walther Janous and Gerhard Kirchner, Innsbruck)
Solution. Throughout the solution we will use the notation xn+1 = x1 .
   We prove that the set of possible values of a is
                                    (                                 )
                                  1 [          1                      n
                             −∞,                     ; k ∈ N, 1 ≤ k <     .
                                  4        4 cos2 kπ
                                                  n
                                                                      2

    In the case a ≤ 14 we can choose x1 such that x1 (1 − x1 ) = a and set x1 = x2 = . . . = xn . Hence we
will now suppose that a > 14 .
    The system (1) gives the recurrence formula
                                                          a    xi − a
                               xi+1 = ϕ(xi ) = 1 −           =        ,         i = 1, . . . , n.
                                                          xi     xi
The fractional linear transform ϕ can be interpreted as a projective transform of the real projective
line R∪ {∞}; the map ϕ is an element of the group PGL2 (R), represented by the linear transform
        1 −a
M =              . (Note that det M 6= 0 since a 6= 0.) The transform ϕn can be represented by M n . A
        1 0
point [u, v] (written in homogenous coordinates) is a fixed point of this transform if and only if (u, v)T
is an eigenvector of M n . Since the entries of M n and the coordinates u, v are real, the corresponding
eigenvalue is real, too.
    The characteristic polynomial of M is x2 − x + a, which has no real root for a > 41 . So M has two
                                              √
conjugate complex eigenvalues λ1.2 = 12 1 ± 4a − 1i . The eigenvalues of M n are λn1,2 , they are real
if and only if arg λ1,2 = ± kπ
                             n
                               with some integer k; this is equivalent with
                                                √               kπ
                                               ± 4a − 1 = tan ,
                                                                 n
                                              1                    1
                                           a=   1 + tan2 kπ
                                                         n
                                                              =           .
                                              4                 4 cos2 kπ
                                                                       n


                                                                   3
   If arg λ1 = kπ  n
                        then λn1 = λn2 , so the eigenvalues of M n are equal. The eigenvalues of M are distinct,
so M and M have two linearly independent eigenvectors. Hence, M n is a multiple of the identity. This
               n

means that the projective transform ϕn is the identity; starting from an arbitrary point x1 ∈ R ∪ {∞},
the cycle x1 , x2 , . . . , xn closes at xn+1 = x1 . There are only finitely many cycles x1 , x2 , . . . , xn containing
the point ∞; all other cycles are solutions for (1).
Remark. If we write xj = P + Q tan tj where P, Q and t1 , . . . , tn are real numbers, the recurrence relation
re-writes as
                                      (P + Q tan tj )(1 − P − Q tan tj+1 ) = a
              (1 − P )Q tan tj − P Q tan tj+1 = a + P (P − 1) + Q2 tan tj tan tj+1   (j = 1, 2, . . . , n).
                                       tan α − tan β                                                     q
In view of the identity tan(α − β) =                  , it is reasonable to choose P = 21 , and Q = a − 41 . Then
                                      1 + tan α tan β
the recurrence leads to                                   √
                                     tj − tj+1 ≡ arctan 4a − 1 (mod π).

Problem 5. Let c ≥ 1 be a real number. Let G be an abelian group and let A ⊂ G be a finite set
satisfying |A + A| ≤ c|A|, where X + Y := {x + y | x ∈ X, y ∈ Y } and |Z| denotes the cardinality of
Z. Prove that
                                                              k
                                     |A
                                      | +A+ {z. . . + A} | ≤ c |A|
                                                    k times
for every positive integer k. (Plünnecke’s inequality)
                                             (Proposed by Przemyslaw Mazur, Jagiellonian University)
Solution. Let B be a nonempty subset of A for which the value of the expression c1 = |A+B||B|
                                                                                               is the
least possible. Note that c1 ≤ c since A is one of the possible choices of B.
Lemma 1. For any finite set D ⊂ G we have |A + B + D| ≤ c1 |B + D|.
Proof. Apply induction on the cardinality of D. For |D| = 1 the Lemma is true by the definition of c1 .
Suppose it is true for some D and let x 6∈ D. Let B1 = {y ∈ B | x + y ∈ B + D}. Then B + (D ∪ {x})
decomposes into the union of two disjoint sets:
                                                                           
                             B + (D ∪ {x}) = (B + D) ∪ (B \ B1 ) + {x}
and therefore |B + (D ∪ {x})| = |B + D| + |B| − |B1 |. Now we need to deal with the cardinality of the
set A + B + (D ∪ {x}). Writing A + B + (D ∪ {x}) = (A + B + D) ∪ (A + B + {x}) we count some
of the elements twice: for example if y ∈ B1 , then A + {y} + {x} ⊂ (A + B + D) ∩ (A + B + {x}).
Therefore all the elements of the set A + B1 + {x} are counted twice and thus
  |A + B + (D ∪ {x})| ≤ |A + B + D| + |A + B + {x}| − |A + B1 + {x}| =
             = |A + B + D| + |A + B| − |A + B1 | ≤ c1 (|B + D| − |B| − |B1 |) = c1 |B + (D ∪ {x})|,
where the last inequality follows from the inductive hypothesis and the observation that |A+B|
                                                                                           |B|
                                                                                               = c1 ≤
|A+B1 |
  |B1 |
        (or B1 is the empty set).                                                                                    2
Lemma 2. For every k ≥ 1 we have | A     . . + A} +B| ≤ ck1 |B|.
                                   | + .{z
                                                k times
Proof. Induction on k. For k = 1 the statement is true by definition of c1 . For greater k set D =
A     . . + A} in the previous lemma: | |A + .{z
| + .{z                                        . . + A} +B| ≤ c1 | A     . . + A} +B| ≤ ck1 |B|.
                                                                   | + .{z                       2
  k−1 times                                    k times                   k−1 times
    Now notice that
                             |A     . . + A} | ≤ | A
                              | + .{z                    . . + A} +B| ≤ ck1 |B| ≤ ck |A|.
                                                   | + .{z
                                  k times                 k times

Remark. The proof above due to Giorgios Petridis and can be found at http://gowers.wordpress.com/
2011/02/10/a-new-way-of-proving-sumset-estimates/

                                                               4
